\documentclass[10pt,a4paper]{amsart}
\usepackage[margin=19mm]{geometry}
\usepackage{amsmath,amssymb,mathtools}
\usepackage[scr=rsfs,cal=euler]{mathalfa}
\usepackage{xfrac}
\usepackage[unicode,hidelinks,bookmarks=false]{hyperref}
\newcommand{\ie}{i.e.\ }
\newcommand{\cf}{cf.\ }
\newcommand{\resp}{respectively\ }
\newcommand{\Subsection}{\subsection}
\providecommand{\nobreakdash}{\nobreak\hbox{-}}
\setlength{\emergencystretch}{2em}
\multlinegap0pt
\usepackage{amssymb}
\usepackage{enumerate}
\newtheorem{theorem}{Theorem}
\newtheorem{lemma}{Lemma}
\title{On sets of subspaces with restricted hyperplane intersection numbers}
\author{Tim Alderson, Simeon Ball}
\date{Source: arXiv:2603.27689v2 (CC BY 4.0)}
\begin{document}
\maketitle

\noindent\emph{Source and licence.} On sets of subspaces with restricted hyperplane intersection numbers. Version arXiv:2603.27689v2 (CC BY 4.0), distributed by arXiv under the Creative Commons Attribution 4.0 International licence, http://creativecommons.org/licenses/by/4.0/, as recorded in the arXiv licence metadata for this exact version and shown on the abstract page https://arxiv.org/abs/2603.27689v2. Full licence notice: \texttt{licenses/arxiv-2603.27689-CC-BY-4.0.txt}.\par\smallskip

\noindent\textit{Editorial scope.} Counted unit: the article's lemma secants (source0.tex lines 223--239). The article's complete original proof of it is delivered with it (source0.tex lines 241--286, one \texttt{proof} environment of 46 lines). No mathematics is written, repaired or re-derived: the statement and the proof are the article's own lines, in the article's own notation, and no retained line is substituted or re-flowed.  Retained uncounted as the source-local context the proof needs: lines 116--126; lines 288--290.  Nothing was omitted from any retained span, so the package carries no omission record.  No retained line contains a citation, so the wrapper carries no bibliography and no bibliography entry.  No retained line contains a figure environment, an image-inclusion command or drawing code, so no image is delivered and the package declares no asset dependency.  Editorial formatting only, in addition: an AMS \texttt{amsart} preamble that restates the article's own declarations and macros used by the retained lines, namely the environments \texttt{theorem}, \texttt{lemma} on separate counters, together with the standard packages those lines need.  The retained environments print in this document as Theorem 1, Lemma 1 in order of appearance, whereas the article prints the same items with the numbering of its own full document.  The complete native source of the article as distributed in the arXiv e-print for this exact version is bundled unchanged as \texttt{sources/197375/source0.tex}.
	\begin{theorem} \label{mainthm}
		Let $\mathcal{X}$ be a multi-set of $(h-1)$-dimensional subspaces in $\mathrm{PG}(kh-1,q)$ with the property that every hyperplane contains at most $t$ subspaces of $\mathcal{X}$. Then
		\begin{equation} \label{eqn: maximal arc bound}
			|\mathcal{X}| \le (t-k+2)q^h+t,
		\end{equation}
		with equality if and only if
		\begin{enumerate}
			\item every set of $(k-1)$ elements of $\mathcal{X}$ spans a $((k-1)h-1)$-dimensional subspace, and \label{part: mainthm 1}
			\item every hyperplane contains either $0, 1, \ldots, k-3$ or $t$ elements of $\mathcal{X}$. \label{part: mainthm 2}
		\end{enumerate}
	\end{theorem}

	\begin{theorem}\label{lem: k=4 maximal arc DNE}
		Let $\mathcal{X}$ be a length-maximal set of $(h-1)$-dimensional subspaces in $\mathrm{PG}(4h-1, q)$ with $q^h > 2$. Then $t=q^h+1$ and every hyperplane is either a $t$-secant or a $1$-secant of $\mathcal{X}$.
	\end{theorem}

	\begin{lemma}\label{lem: enumerating secants}
		Let $\mathcal{X}$ be a length-maximal set of $(h-1)$-dimensional subspaces in $\mathrm{PG}(kh-1,q)$. Set $n = |\mathcal{X}| = (t-k+2)q^h+t$. Then
		
		\begin{enumerate}
			\item $\displaystyle N(t,h,3,t) = \frac{q^{3h}-1}{q-1} -\frac{q^h-1}{q-1}\cdot\left(\frac{q^h(q^h+1)}{t}-q^h\right)$ \label{part 1: lem: enumerating secants}
			
			\item $\displaystyle N(t,h,3,0) = \frac{q^h-1}{q-1}\cdot\left(\frac{q^h(q^h+1)}{t}-q^h\right)$ \label{part 2: lem: enumerating secants}
			
			\item $\displaystyle N(t,h,4,t) = \left( q^h + 1 - \frac{2q^h}{t} \right) \left( q^h + 1 - \frac{q^h}{t-1} \right) \frac{q^{2h}-1}{q-1}$ \label{part 3: lem: enumerating secants}
			
			\item $\displaystyle N(t,h,4,1) = n\cdot N(t-1,h,3,0) =\left((t-2)q^h + t\right) \cdot \frac{q^h-1}{q-1}\cdot \left(  \frac{q^h(q^h+1)}{t-1}-q^h \right)$ \label{part 4: lem: enumerating secants}
			
			\item $\displaystyle N(q^h+2,h,5,2) = \frac{n}{2}\cdot N(q^h+1,h,4,1) = \frac{q^{2h}+2}{2}\cdot (q^{2h}+1)\cdot\frac{q^h-1}{q-1}$ \label{part 5: lem: enumerating secants}
			
			\item $\displaystyle N(q^h+2,h,5,q^h+2) = \frac{(q^{2h}+2)(q^{2h}+1)q^h(q^h-1)}{(q^h+2)(q-1)}$ \label{part 6: lem: enumerating secants}
		\end{enumerate}
	\end{lemma}

	\begin{proof}
		Set $n = (t-k+2)q^h+t$ and we use the term $(h-1)$-flat for an $(h-1)$-dimensional subspace.
		
		\medskip\noindent\textit{Part~\ref{part 1: lem: enumerating secants}.}
		By part~\ref{part: mainthm 2} of Theorem~\ref{mainthm}, every hyperplane is a $0$-secant or a $t$-secant when $k=3$. Counting pairs $(x,H)$ with $x \in \mathcal{X}$ and $H$ a $t$-secant through $x$, and using the fact that each $(h-1)$-flat in $\mathrm{PG}(3h-1,q)$ lies in $\frac{q^{2h}-1}{q-1}$ hyperplanes,
		$$
		N(t,h,3,t)\cdot t = n\cdot \frac{q^{2h}-1}{q-1}.
		$$
		Substituting $n=(t-1)q^h+t$ and rearranging gives the formula. Part~\ref{part 2: lem: enumerating secants} follows since $N(t,h,3,0)+N(t,h,3,t)=|\mathrm{PG}(3h-1,q)|$.
		
		\medskip\noindent\textit{Part~\ref{part 3: lem: enumerating secants}.}
		By part~\ref{part: mainthm 1} of Theorem~\ref{mainthm}, any two elements of $\mathcal{X}$ span a $(2h-1)$-flat; each such flat lies in $\frac{q^{2h}-1}{q-1}$ hyperplanes of $\mathrm{PG}(4h-1,q)$. By part~\ref{part: mainthm 2} of Theorem~\ref{mainthm}, each hyperplane is a $0$-, $1$-, or $t$-secant when $k=4$. Counting triples $(x,y,H)$ with $x \ne y$ in $\mathcal{X}$ and $H$ a $t$-secant through both,
		$$
		N(t,h,4,t)\cdot \binom{t}{2} = \binom{n}{2} \cdot \frac{q^{2h}-1}{q-1}.
		$$
		Substituting $n=(t-2)q^h+t$ and simplifying gives the formula.
		
		\medskip\noindent\textit{Part~\ref{part 4: lem: enumerating secants}.}
		For $k=4$, a $1$-secant hyperplane contains exactly one element $x_i \in \mathcal{X}$. The quotient of $\mathcal{X}$ at any $x_i \in \mathcal{X}$ is a length-maximal set $\mathcal{X}_i^*$ in $\mathrm{PG}(3h-1,q)$ with parameter $t-1$. Each $0$-secant of $\mathcal{X}_i^*$ corresponds to a $1$-secant of $\mathcal{X}$ through $x_i$, so summing over all $x_i$,
		$$
		N(t,h,4,1) = n \cdot N(t-1,h,3,0).
		$$
		Substituting the formula from part~\ref{part 2: lem: enumerating secants} of Lemma~\ref{lem: enumerating secants} with $t$ replaced by $t-1$ gives the result.
		
		\medskip\noindent\textit{Part~\ref{part 5: lem: enumerating secants}.}
		Now let $k=5$ and $t=q^h+2$, so $n=q^{2h}+2$. By part~\ref{part: mainthm 2} of Theorem~\ref{mainthm}, every hyperplane is a $0$-, $1$-, $2$-, or $t$-secant. Since any hyperplane containing $3$ or more elements of $\mathcal{X}$ must contain $t$ elements (as any $3$ elements span a $(3h-1)$-flat, and any hyperplane through a $(3h-1)$-flat that contains a fourth element of $\mathcal{X}$ must contain all $t$ by part~\ref{part: mainthm 2} of Theorem~\ref{mainthm}), the secant types are restricted to $0$, $1$, $2$, and $t$.
		
		For $k=5$, each element $x_i \in \mathcal{X}$ lies in a $2$-secant hyperplane $H$ if and only if $H$ contains exactly one further element $x_j \ne x_i$. The number of $2$-secant hyperplanes through $x_i$ equals the number of $1$-secant hyperplanes of the quotient set $\mathcal{X}_i^*$ in $\mathrm{PG}(4h-1,q)$, which has parameter $t-1=q^h+1$. Summing over all $x_i$ and dividing by $2$ (since each $2$-secant is counted twice),
		$$
		N(q^h+2,h,5,2) = \frac{n}{2}\cdot N(q^h+1,h,4,1).
		$$
		By Theorem~\ref{lem: k=4 maximal arc DNE} (proved below), the quotient set $\mathcal{X}_i^*$ has $t-1=q^h+1$ and has no $0$-secants, so $N(q^h+1,h,4,1) = |\mathrm{PG}(4h-1,q)| - N(q^h+1,h,4,q^h+1)$. Alternatively, substituting $t-1=q^h+1$ directly into part~\ref{part 4: lem: enumerating secants} of Lemma~\ref{lem: enumerating secants} and simplifying gives
		$$
		N(q^h+2,h,5,2) = \frac{q^{2h}+2}{2}\cdot (q^{2h}+1)\cdot\frac{q^h-1}{q-1}.
		$$
		
		\medskip\noindent\textit{Part~\ref{part 6: lem: enumerating secants}.}
		With $k=5$, $t=q^h+2$, $n=q^{2h}+2$, any three elements of $\mathcal{X}$ span a $(3h-1)$-flat (since any $k-1=4$ elements span a $(4h-1)$-flat by part~\ref{part: mainthm 1} of Theorem~\ref{mainthm}, which implies any three also span a $(3h-1)$-flat). Each $(3h-1)$-flat in $\mathrm{PG}(5h-1,q)$ lies in $\frac{q^{2h}-1}{q-1}$ hyperplanes. Counting quadruples $(x_1,x_2,x_3,H)$ with distinct $x_i \in \mathcal{X}$ and $H$ a $t$-secant through all three,
		$$
		N(q^h+2,h,5,q^h+2)\cdot \binom{t}{3} = \binom{n}{3}\cdot\frac{q^{2h}-1}{q-1}.
		$$
		Substituting $n=q^{2h}+2$ and $t=q^h+2$ and simplifying,
		$$
		N(q^h+2,h,5,q^h+2) = \frac{(q^{2h}+2)(q^{2h}+1)q^{2h}}{(q^h+2)(q^h+1)q^h}\cdot\frac{q^{2h}-1}{q-1} = \frac{(q^{2h}+2)(q^{2h}+1)q^h(q^h-1)}{(q^h+2)(q-1)}.
		$$
	\end{proof}

\end{document}
